\documentclass[11pt]{llmxive}

\title{A Stylometric Application of Large Language Models}
\author{Harrison F. Stropkay \and Jiayi Chen \and Mohammad J. Latifi \and
        Daniel N. Rockmore \and Jeremy R. Manning\\[0.3em]
        \normalsize\itshape Dartmouth College, Hanover, NH 03755, USA}
\date{May 2026}
\arxivid{2510.21958}
\githubrepo{https://github.com/ContextLab/llm-stylometry}

\begin{document}
\maketitle

\begin{abstract}
We show that large language models (LLMs) can be used to distinguish the
writings of different authors. Specifically, an individual GPT-2 model,
trained from scratch on the works of one author, will predict held-out text
from that author more accurately than held-out text from other authors. We
suggest that, in this way, a model trained on one author's works embodies the
unique writing style of that author. We first demonstrate our approach on
books written by eight different (known) authors. We also use this approach to
confirm R. P. Thompson's authorship of the well-studied 15th book of the Oz
series, originally attributed to F. L. Baum.
\end{abstract}

\section{Introduction}
Herein we introduce \emph{predictive comparison}, a new LLM-based relative
stylometric measure. It derives from a simple idea, that if an LLM can be
trained to write like a given author by training on their work, then the
degree to which such a model can predict held-out text from that author
should reflect the strength of the author's style.

Stylometry has a long history in computational linguistics. Classical work
relies on lexical, syntactic, and structural features---word frequencies,
function-word distributions, part-of-speech \emph{n}-grams. Modern neural
methods replace hand-crafted features with learned representations.

\section{Approach}
For each candidate author $A$ in a set, we train a small GPT-2 model on a
training partition of $A$'s writings, then evaluate cross-entropy on
held-out passages from every author in the set.

\begin{itemize}
  \item \textbf{Content tokens:} the full token stream.
  \item \textbf{Function words:} content words masked out.
  \item \textbf{POS tags:} every token replaced by its part-of-speech.
\end{itemize}

We use Python and PyTorch; the full training and evaluation code is at
\href{https://github.com/ContextLab/llm-stylometry}{ContextLab/llm-stylometry}.

\subsection{Cross-entropy}
For an author $A$ trained model $\theta_A$ and a held-out passage $x$ from
author $B$, the per-token cross-entropy is
\[
  H(x;\theta_A) = -\tfrac{1}{|x|}\sum_{t=1}^{|x|} \log p_{\theta_A}(x_t \mid x_{<t}).
\]
The lower the cross-entropy, the more author-$B$-like the model finds the
passage.

\subsubsection{Statistical comparison}
We compare cross-entropies via two-sample \emph{t}-tests over passages.

\section*{Discussion}
A small accent rule, a colored section number, a few links: that's what we
need to look at to verify the house style is doing what we want.

\section*{Code availability}
Code: \url{https://github.com/ContextLab/llm-stylometry}.

\end{document}
